\subsubsection*{Informal Proof Patterns}

\begin{remark*}[Working vocabulary]
This topic connects valid argument forms with common informal proof patterns:
direct proof, conditional proof, proof by cases, contradiction, contrapositive
proof, implication chains, and equivalence chains.
\end{remark*}

% BEGIN GENERATED ARTIFACT: def:informal-proof-pattern-propositional-logic
\begin{definitionbox}{Definition (Informal Proof Pattern)}
\begin{definition}[Informal Proof Pattern]
\label{def:informal-proof-pattern-propositional-logic}
An \emph{informal proof pattern} is a reusable truth-preserving strategy for
organizing an ordinary mathematical proof.
\end{definition}
\end{definitionbox}

\begin{remark*}[Standard quantified statement]
\[
\Pi\in\mathsf{InformalProofPattern}
\Longrightarrow
\Pi\text{ is justified by one or more valid argument forms.}
\]
\end{remark*}

\begin{remark*}[Predicate reading]
\(\Pi\in\mathsf{InformalProofPattern}\) means that \(\Pi\) is a recognized
proof strategy whose correctness rests on truth-preserving inference.
\end{remark*}

\begin{remark*}[Interpretation]
Informal proof patterns are not a substitute for proof. They are templates for
arranging valid moves so that a proof has a clear logical shape.
\end{remark*}

\begin{remark*}[Exposition]
A direct proof establishes a target conclusion by moving forward from known
premises through valid argument forms. A conditional proof establishes
\(\varphi\to\psi\) by assuming \(\varphi\) and deriving \(\psi\) inside that
assumption context. Proof by cases uses a disjunction of cases and shows that
each case leads to the desired conclusion.
\end{remark*}

\begin{remark*}[Exposition]
Proof by contradiction assumes the negation of the desired conclusion and derives
an impossible situation. Contrapositive proof establishes \(\varphi\to\psi\) by
showing \(\neg\psi\to\neg\varphi\). Chain-of-implications proofs use repeated
hypothetical syllogism, while equivalence-chain proofs use replacement and
logical equivalence.
\end{remark*}

\begin{remark*}[Exposition]
If the target has the form \(\varphi\to\psi\), a conditional proof is often
natural. If the available information has the form \(\varphi\lor\psi\), proof by
cases is often natural. If the target is an equivalence, an equivalence chain or
two conditional proofs are often natural.
\end{remark*}

\begin{dependencies}
\begin{itemize}
  \item \hyperref[def:valid-argument-form-propositional-logic]{Valid Argument Form}
  \item \hyperref[def:logical-consequence-propositional-logic]{Logical Consequence}
\end{itemize}
\end{dependencies}
% END GENERATED ARTIFACT: def:informal-proof-pattern-propositional-logic
